\documentclass[12pt,a4paper]{ctexart}

% 页面设置
\usepackage[a4paper,margin=2.5cm]{geometry}
\usepackage{setspace}
\usepackage{titling}

\begin{document}

% ==================== 封面 ====================
\begin{titlepage}
    \centering

    % 项目名称
    \vspace*{3.5cm}

    {\zihao{1}\bfseries 智序工作台\par}

    \vspace{0.5cm}

    {\zihao{3}\bfseries （Zhixu Workbench）\par}

    \vspace{1.5cm}

    {\zihao{1}\bfseries 创业计划书\par}

    \vfill

    % 团队信息
    \begin{flushleft}
        \hspace{2.5cm}
        \begin{minipage}{0.7\textwidth}
            \zihao{4}
            \setlength{\baselineskip}{1.8em}

            \textbf{学校：南京林业大学}

            \textbf{班级：25108011}

            \textbf{团队成员：xxx}

            \hspace*{3.2em} xxx

            \hspace*{3.2em} xxx
        \end{minipage}
    \end{flushleft}

    \vspace{2cm}

    % 日期
    {\zihao{4}\bfseries 2026 年 10 月\par}

    \vspace{1.5cm}

\end{titlepage}

% ==================== 正文开始 ====================
\tableofcontents
\newpage

\section{项目概述}

在此填写智序工作台项目的简介、创业背景及核心价值。

\section{市场分析}

在此填写市场需求、目标用户及竞争环境分析。

\section{产品与服务}

在此填写智序工作台的产品功能、技术方案及服务模式。

\section{商业模式}

在此填写项目的盈利方式、收入来源及运营策略。

\section{团队介绍}

在此填写团队成员、职责分工及相关经验。

\section{财务分析}

在此填写项目预算、成本估算及收益预测。

\section{风险分析与应对措施}

在此填写项目可能面临的风险及对应的解决方案。

\end{document}
